\documentclass[conference]{IEEEtran}

\usepackage{graphicx}
\usepackage{amsmath,amssymb}
\usepackage{float}
\usepackage{booktabs}
\usepackage{multirow}

\begin{document}

\title{
    \textbf{Benchmark Report: 510.parest\_r} \\
    \text{Register Spill Analysis on RISC-V via gem5}
}

\author{
    \IEEEauthorblockN{Lütfullah Burak Kaya}
    \IEEEauthorblockA{
        Ozyegin University \\
        burak.kaya.50711@ozu.edu.tr
    }
    \and
    \IEEEauthorblockN{Ismail Akturk}
    \IEEEauthorblockA{
        Ozyegin University \\
        ismail.akturk@ozyegin.edu.tr
    }
}

\newcommand{\LongDate}{%
    \number\day\space
    \ifcase\month
    \or January\or February\or March\or April\or May\or June%
    \or July\or August\or September\or October\or November\or December%
    \fi
    \space \number\year
}

\twocolumn[
\begin{@twocolumnfalse}
\maketitle
\begin{center}{\small \LongDate}\end{center}
\vspace{1em}
\end{@twocolumnfalse}
]

% ================================================

\begin{abstract}
This report presents the register spill characterization of the
\texttt{510.parest\_r} benchmark from SPEC CPU2017, executed on a
RISC-V (RV64GC) target using the gem5 TimingSimpleCPU simulator.
A custom SpillDetector was integrated into gem5 to count store-then-load
patterns on the stack within a user-defined Region of Interest (ROI).
The test dataset (\texttt{test.prm}, 8$\times$8$\times$4 coarse mesh,
1 experiment) was used as the input. Results show a spill rate of
\textbf{0.34\%} within the ROI --- the lowest of all 17 evaluated
SPEC CPU2017 benchmarks --- with 128 million detected spill events
across 37.3 billion ROI instructions. The parest finite element solver
is dominated by sparse matrix--vector products and conjugate gradient
iterations, which exhibit a highly regular memory access pattern and
an extreme load-to-store ratio of 11.8:1. The register file is
effectively utilized by the iterative solver kernels, leaving
almost no spill pressure on the RISC-V register allocator.
\end{abstract}

% =========================================================
\section{Benchmark Overview}

\texttt{510.parest\_r} is the SPEC CPU2017 rate variant of
\emph{par}allel \emph{est}imation (parest), a C++ application for
parameter estimation in partial differential equations, specifically
medical diffuse optical tomography. The benchmark solves an
inverse problem: given synthetic boundary measurements of light
transport in a 3-D tissue phantom, recover the spatially varying
absorption coefficient using an iterative Newton-CG method backed
by the deal.II finite element library.

The computational core involves:
\begin{itemize}
    \item Assembling the finite element stiffness matrix on a
          3-D hexahedral mesh (8$\times$8$\times$4 elements for the
          test dataset);
    \item Solving forward and adjoint PDEs via conjugate gradient
          with AMG preconditioning;
    \item Evaluating the misfit functional and its gradient;
    \item Performing mesh refinement and updating the parameter
          discretisation across Newton iterations.
\end{itemize}

The RISC-V binary (\texttt{parest\_r.riscv}) was compiled with gem5
ROI markers inserted around the \texttt{do\_job()} call in
\texttt{src/source/me-tomography/me\_tomography.cc}:

\begin{itemize}
    \item \texttt{m5\_work\_begin(0,0)} --- before \texttt{do\_job()}
    \item \texttt{m5\_work\_end(0,0)}   --- after  \texttt{do\_job()}
\end{itemize}

\texttt{do\_job()} encompasses all Newton iterations including mesh
assembly, forward/adjoint solves, gradient evaluation, and adaptive
mesh refinement.

\subsection{Input Datasets}

SPEC CPU2017 provides three dataset sizes for \texttt{parest\_r},
differing in mesh resolution and number of Newton iterations
(Table~\ref{tab:inputs}).

\begin{table}[H]
\centering
\caption{parest\_r Input Dataset Sizes}
\label{tab:inputs}
\begin{tabular}{lll}
\toprule
\textbf{Dataset} & \textbf{Mesh} & \textbf{Description} \\
\midrule
test    & 8$\times$8$\times$4    & 1 experiment, coarse mesh \\
train   & medium                 & more refinement levels \\
refrate & fine                   & full production run \\
\bottomrule
\end{tabular}
\end{table}

The \textbf{test} dataset was selected for this study. Despite the
coarse mesh, the Newton-CG solver generates 37.3 billion RISC-V
instructions because each conjugate gradient iteration on even a
small 3-D FEM mesh requires numerous sparse matrix--vector products,
each traversing large irregular data structures in memory.

\subsection{gem5 SE Mode Compatibility}

The parest source accesses \texttt{/proc/loadavg} to report system
load in diagnostic output. This virtual file is not available in
gem5's Syscall Emulation (SE) mode. The single call site in
\texttt{src/source/base/utilities.cc} (\texttt{get\_cpu\_load()})
was patched to return 0.0 unconditionally; this value is used
only for logging and has no effect on the simulation result.

% =========================================================
\section{Simulation Setup}

\subsection{gem5 Configuration}

\begin{itemize}
    \item \textbf{CPU model:}   TimingSimpleCPU
    \item \textbf{ISA:}         RISC-V RV64GC
    \item \textbf{Caches:}      enabled (default L1 I/D + L2)
    \item \textbf{Memory:}      4\,GB simulated physical RAM
    \item \textbf{Spill mode:}  summary only (\texttt{spill\_verbose=False})
\end{itemize}

The full gem5 invocation was:

\begin{verbatim}
./build/RISCV/gem5.opt
  --outdir=benchmarks/parest/m5out_test
  configs/deprecated/example/se.py
  --cmd=benchmarks/parest/parest_r.riscv
  --options="benchmarks/parest/test.prm"
  --cpu-type=TimingSimpleCPU
  --caches
  --mem-size=4GB
\end{verbatim}

The single argument \texttt{test.prm} is the deal.II parameter file
that specifies the mesh dimensions, FEM discretisation, Newton solver
tolerances, and output format.

% =========================================================
\section{Simulation Results}

The simulation completed successfully. The Newton solver executed all
configured iterations, produced VTK output files for the solution,
and exited cleanly.

\subsection{Pre-ROI Context}

Before \texttt{do\_job()} was entered, the binary executed deal.II
runtime initialisation: parsing the parameter file, constructing the
finite element and quadrature objects, and building the coarse mesh.
Table~\ref{tab:preroi} shows the global counters at ROI entry.

\begin{table}[H]
\centering
\caption{Global Counters at ROI Entry}
\label{tab:preroi}
\begin{tabular}{lr}
\toprule
\textbf{Metric} & \textbf{Value} \\
\midrule
Total instructions (pre-ROI) & 4,282,420 \\
Total spills detected (pre-ROI) & 145,908 \\
Pre-ROI spill rate & 3.41\% \\
\bottomrule
\end{tabular}
\end{table}

The pre-ROI spill rate (3.41\%) is an order of magnitude higher than
the ROI rate (0.34\%). This contrast arises because the initialisation
phase exercises deep C++ constructor chains from deal.II (cell
iterators, FE objects, triangulation structures), which generate
significant register pressure, whereas the ROI solver kernels are
dominated by numerically flat inner loops.

\subsection{ROI Spill Statistics}

Table~\ref{tab:roi} presents the spill statistics collected within
the ROI.

\begin{table}[H]
\centering
\caption{ROI Spill Statistics --- 510.parest\_r (test.prm)}
\label{tab:roi}
\begin{tabular}{lr}
\toprule
\textbf{Metric} & \textbf{Value} \\
\midrule
ROI instructions         & 37,262,401,282 \\
ROI stores               &  1,009,678,754 \\
ROI loads                & 11,899,900,782 \\
\midrule
\textbf{ROI spills}      & \textbf{127,688,196} \\
\textbf{Spill rate}      & \textbf{0.3427\%} \\
\bottomrule
\end{tabular}
\end{table}

\subsection{Timing}

\begin{table}[H]
\centering
\caption{Simulation Timing}
\label{tab:timing}
\begin{tabular}{lr}
\toprule
\textbf{Metric} & \textbf{Value} \\
\midrule
Simulated time           & 94.18\,s \\
Host wall time           & 23{,}977.00\,s ($\approx$6.66\,hours) \\
Total simulated instructions & 37,270,019,516 \\
\bottomrule
\end{tabular}
\end{table}

% =========================================================
\section{Analysis}

\subsection{Spill Rate Interpretation}

The measured spill rate of \textbf{0.34\%} places
\texttt{parest\_r} seventeenth (last) among all 17 evaluated SPEC
CPU2017 benchmarks (Table~\ref{tab:compare}), well below the
previous minimum of 1.03\% (\texttt{508.namd\_r}).

\begin{table}[H]
\centering
\caption{Cross-Benchmark Spill Rate Comparison}
\label{tab:compare}
\begin{tabular}{lrr}
\toprule
\textbf{Benchmark} & \textbf{ROI Spills} & \textbf{Spill Rate} \\
\midrule
507.cactuBSSN\_r & 3,782,099,438 & 7.8756\% \\
511.povray\_r    &   174,359,922 & 7.7809\% \\
500.perlbench\_r &     1,070,055 & 7.2838\% \\
520.omnetpp\_r   &   894,619,681 & 7.1994\% \\
502.gcc\_r       &     1,108,838 & 6.9185\% \\
526.blender\_r   &    61,845,869 & 6.2193\% \\
523.xalancbmk\_r &    18,381,250 & 4.6427\% \\
541.leela\_r     & 1,158,236,271 & 4.4438\% \\
519.lbm\_r       &   274,130,440 & 4.2130\% \\
531.deepsjeng\_r & 1,928,903,656 & 4.1749\% \\
525.x264\_r      &   767,981,981 & 3.2812\% \\
544.nab\_r       &   151,310,616 & 2.2612\% \\
538.imagick\_r   &     2,424,899 & 2.0601\% \\
557.xz\_r        &    38,411,428 & 1.9069\% \\
505.mcf\_r       &   446,113,136 & 1.8263\% \\
508.namd\_r      &   327,718,955 & 1.0269\% \\
510.parest\_r    &   127,688,196 & \textbf{0.3427\%} \\
\bottomrule
\end{tabular}
\end{table}

\subsection{Store vs.\ Load Balance}

The load count (11.90B) is $\mathbf{11.79\times}$ higher than the
store count (1.01B), the most extreme load dominance observed across
all 17 evaluated benchmarks. This structural characteristic of the
finite element solver arises directly from the conjugate gradient
algorithm:
\begin{itemize}
    \item \textbf{Sparse matrix--vector products (SpMV)}: each SpMV
          traverses the entire non-zero structure of the stiffness
          matrix (loads from \texttt{values[]} and \texttt{col\_ind[]}
          arrays) and accumulates results into a dense output vector
          (one store per row). The non-zero count greatly exceeds the
          row count, producing many loads per store.
    \item \textbf{Vector updates} (\texttt{axpy}, dot products): these
          read two input vectors and write one output, yielding 2:1
          load:store ratios that compound across the CG loop.
    \item \textbf{AMG preconditioner application}: similar SpMV
          pattern at each level of the multigrid hierarchy.
\end{itemize}
Across all CG iterations, the cumulative effect is an 11.79:1
load-to-store ratio, far exceeding the 2--3:1 typical of
object-oriented or tree-traversal workloads.

\subsection{Spill Fraction of Loads}

Within the ROI, the ratio of spills to total loads is:
\[
\frac{127{,}688{,}196}{11{,}899{,}900{,}782} \approx 1.07\%
\]
Fewer than 1 in 90 load instructions is a spill reload. The FEM
solver inner loop operates on a small set of variables: the current
residual vector pointer, the search direction pointer, two scalar
step-size values ($\alpha$, $\beta$), and the iteration counter.
RISC-V's 32 registers accommodate this working set trivially, leaving
virtually no register pressure. The 1.07\% spill-to-load fraction
is the lowest of all evaluated benchmarks, reflecting the
register-friendly nature of iterative Krylov solvers.

\subsection{Pre-ROI vs.\ ROI Spill Rate Contrast}

The pre-ROI spill rate (3.41\%) is $10\times$ higher than the ROI
rate (0.34\%). This inversion is structurally caused by the layered
initialisation of the deal.II library: constructing a
\texttt{Triangulation}, a \texttt{DoFHandler}, and a
\texttt{ConstraintMatrix} involves deep C++ constructor chains with
many arguments passed by reference across multiple abstraction layers,
generating register pressure comparable to object-oriented benchmarks.
Once initialisation completes and the solver enters the Newton loop,
the computational character shifts entirely to numerical kernels with
minimal function-call depth.

\subsection{ROI Coverage}

The ROI accounts for $37{,}262{,}401{,}282$ out of $37{,}270{,}019{,}516$
total simulated instructions:
\[
\frac{37{,}262{,}401{,}282}{37{,}270{,}019{,}516} \approx 99.98\%
\]
The near-perfect ROI coverage confirms that \texttt{do\_job()}
encompasses virtually all benchmark computation. The 0.02\% outside
the ROI corresponds to deal.II mesh construction and parameter file
parsing.

\subsection{FEM Solver Register Efficiency}

The 0.34\% spill rate demonstrates that iterative finite element
solvers are among the most register-efficient workloads on RISC-V
RV64GC. The CG and SpMV kernels that dominate execution share two
structural properties:
\begin{itemize}
    \item \textbf{Flat call depth}: the hot path is a tight loop with
          no virtual dispatch or recursive calls, so the compiler can
          keep all live variables in registers without spilling.
    \item \textbf{Memory bandwidth bound}: the bottleneck is cache
          miss latency during sparse matrix traversal, not
          register-file capacity. Spills would add additional memory
          traffic that the compiler actively avoids.
\end{itemize}
The extreme 11.79:1 load-to-store ratio thus reflects the memory
access pattern of the algorithm, not register pressure: the register
file is not oversubscribed, but the working data structures (sparse
matrices, multi-level vectors) far exceed cache capacity and are
streamed from DRAM on every iteration.

% =========================================================
\section{Conclusion}

The \texttt{510.parest\_r} benchmark exhibits a register spill rate
of 0.34\% within the \texttt{do\_job()} ROI on RISC-V RV64GC, the
lowest of all 17 evaluated SPEC CPU2017 benchmarks. With 128 million
spill events across 37.3 billion ROI instructions and a 99.98\%
ROI coverage, the benchmark is comprehensively characterised. The
11.79:1 load-to-store ratio is the highest of the entire benchmark
suite, reflecting the streaming sparse matrix--vector access pattern
of the underlying conjugate gradient solver. The 1.07\% spill-to-load
fraction and the sharp 10$\times$ contrast between pre-ROI (3.41\%)
and ROI (0.34\%) spill rates together confirm that iterative FEM
solvers are naturally register-efficient on RISC-V: the sparse
numerical kernels that dominate execution keep a minimal working set
in registers, while the algorithmic bottleneck is memory bandwidth
rather than register-file capacity.

\end{document}
